\documentclass[11pt,a4paper]{article}
\usepackage[T1]{fontenc}
\usepackage[utf8]{inputenc}
\usepackage{lmodern,amsmath,amssymb}
\usepackage[margin=25mm]{geometry}
\usepackage{longtable,booktabs,array,calc}
\usepackage[hidelinks]{hyperref}
\hypersetup{pdftitle={In three dimensions the mass gap is the positivity of one function at infinity},pdfauthor={navstokgap research working notes}}
\newcommand{\tightlist}{\setlength{\itemsep}{0pt}\setlength{\parskip}{0pt}}
\setlength{\emergencystretch}{3em}
\setcounter{secnumdepth}{0}
\begin{document}
\hypertarget{in-three-dimensions-the-mass-gap-is-the-positivity-of-one-function-at-infinity}{%
\section{In three dimensions the mass gap is the positivity of one
function at
infinity}\label{in-three-dimensions-the-mass-gap-is-the-positivity-of-one-function-at-infinity}}

For pure Yang--Mills in \(d=3\) (two space dimensions) the coupling
\(g^2\) is itself a mass, so the finite-volume gap of the continuum
theory on a periodic square torus of side \(L\) has the exact form
\[\Delta(L)=g^2\hbar^2c\,f(x),\qquad x=g^2\hbar L\ \text{(dimensionless)},\]
by dimensional analysis alone, and the Jaffe--Witten conjecture for
\(d=3\) is the single statement
\[\lim_{x\to\infty}f(x)=C\in(0,\infty),\qquad m=C\,g^2\hbar^2/c .\] The
small-\(x\) end of \(f\) is known from the constant-mode sector: with
two spatial components the zero-mode Hamiltonian is C133 with \(D=2\),
\(m_{\rm eff}=L^2/(g^2c)\) and \(g_B^2=L^2c/g^2\), giving
\(\Delta_0(L)=\delta_1^{(2)}\,g^2\hbar^2c\,x^{-2/3}\), that is
\(f(x)=\delta_1^{(2)}x^{-2/3}[1+o(1)]\) as \(x\to0\). On the lattice,
the strong-coupling theorem of \href{strong-coupling-uniform-gap.md}{the
previous note} with \(\nu=2\) gives a gap
\(\ge\gamma_2(g_{\rm lat}^2/2)C_2\hbar c/a\) uniform in the lattice
size, and the continuum limit is the statement that
\(\delta_\infty(g_{\rm lat})/g_{\rm lat}^2\) stays bounded below as
\(g_{\rm lat}^2=g^2\hbar a\to0\): the gap-to-coupling ratio is bounded
below at both ends of the coupling axis, and the conjecture is its
positivity in between. Contrast \(d=4\), where the ratio must pass from
a power law \(\propto g^2\) at strong coupling to the essential
singularity \(e^{-1/(2b_0g^2)}\) at weak coupling. The three-dimensional
problem is therefore the natural first target: it has no transmutation,
a single scaling variable, and both ends of \(f\) in hand. The expected
value \(C=N/(2\pi)\) (Karabali--Nair, physics argument; metadata, B78)
is a prediction to be proved, not an input. Nothing here is promoted.

\hypertarget{units-and-the-one-variable-form}{%
\subsection{1. Units and the one-variable
form}\label{units-and-the-one-variable-form}}

With \(S=\frac1{4g^2}\int F^a_{\mu\nu}F^{a\mu\nu}\,d^3x\) and the
quantization \(e^{iS/\hbar}\), the dimensional census of
\href{low-dimensional-mass-gap.md}{G07} §4 gives \([1/g^2]=\)
action\(\cdot\)length in \(d=3\), so \(g^2\hbar=1/\ell_g\) is an inverse
length, \(\mu:=g^2\hbar^2/c=\hbar/(c\ell_g)\) is a mass, and
\(\varepsilon_3:=\mu c^2=g^2\hbar^2c\) is the unique energy formed from
\((\hbar,c,g)\). The continuum theory on the periodic torus
\(T^2_L\times\mathbb R\) has the fixed constants \(\hbar\), \(c\), \(g\)
and the box size \(L\), and a gap \(\Delta(L)\), when the finite-volume
theory exists, depends on them only. By the floor/unit theorem C131
applied to the observable ``energy'', \(\Delta(L)\) is \(\varepsilon_3\)
times a function of the single dimensionless combination
\(x=L/\ell_g=g^2\hbar L\), because \((\hbar,c,g)\) span exactly one
energy and \(L\) adds exactly one dimensionless ratio. Setting
\(f(x)=\Delta(L)/\varepsilon_3\):
\[\Delta(L)=\varepsilon_3\,f(x),\qquad\varepsilon_3=g^2\hbar^2c,\qquad x=g^2\hbar L .\]
The infinite-volume gap, if it exists, is
\(m c^2=\varepsilon_3\lim_{x\to\infty}f(x)\), and Jaffe--Witten's
\(0<m<\infty\) is \(0<\lim f<\infty\). Vacuum-energy renormalization,
the only ultraviolet divergence of the superrenormalizable theory
relevant here, shifts \(E_0\) and \(E_1\) equally and leaves \(f\)
untouched.

\hypertarget{the-small-x-end-from-the-constant-modes}{%
\subsection{\texorpdfstring{2. The small-\(x\) end from the constant
modes}{2. The small-x end from the constant modes}}\label{the-small-x-end-from-the-constant-modes}}

Temporal gauge on \(T^2_L\) leaves two constant spatial components
\(a_1,a_2\) in the Lie algebra. With
\(S=\frac1{4g^2}\int F^a_{\mu\nu}F^{a\mu\nu}d^2x\,c\,dt\),
\(F_{0i}=c^{-1}\dot a_i\) and \(F_{12}=\vec a_1\times\vec a_2\), the
constant-mode Lagrangian is
\(\frac{L^2}{2g^2c}|\dot{\vec a}|^2-\frac{L^2c}{2g^2}|\vec a_1\times\vec a_2|^2\)
and
\[H_0=\frac{g^2c}{2L^2}\sum_{i=1}^2|\vec p_i|^2+\frac{L^2c}{2g^2}\,|\vec a_1\times\vec a_2|^2,
\qquad[a_i^a,p_j^b]=i\hbar\delta_{ij}\delta^{ab},\] which is C133 with
\(D=2\), \(m_{\rm eff}=L^2/(g^2c)\) and \(g_B^2=L^2c/g^2\); \(\hbar\)
enters only through \(g^2\) and the commutator. Its energy unit is
\[\varepsilon=\hbar^{4/3}g_B^{2/3}m_{\rm eff}^{-2/3}
=\hbar^{4/3}\Big(\frac{L^2c}{g^2}\Big)^{1/3}\Big(\frac{L^2}{g^2c}\Big)^{-2/3}
=\hbar^{4/3}c\,g^{2/3}L^{-2/3}
=\frac{\hbar c}{L}\Big(\frac{L}{\ell_g}\Big)^{1/3}
=\varepsilon_3\,x^{-2/3},\] using \(g^{2/3}\hbar^{1/3}=\ell_g^{-1/3}\)
and \(\hbar c/L=\varepsilon_3/x\). So the constant-mode gap is
\(\Delta_0(L)=\delta_1^{(2)}\varepsilon_3x^{-2/3}\) with
\(\delta_1^{(2)}>0\) the pure number of C133 for \(D=2\), and
\[f(x)=\delta_1^{(2)}\,x^{-2/3}\,[1+o(1)],\qquad x\to0,\] where \(o(1)\)
stands for the coupling to the nonzero momentum modes, controlled
perturbatively by the small-volume expansion in the same way as in
\(d=4\) (Lüscher's method; the \(d=3\) expansion parameter is
\(x^{1/3}\) rather than \(g^{2/3}\), since \(x\) is the only parameter).
The claim \(f\to\infty\) as \(x\to0\) is the statement that the small
box is gapped with a gap growing like \(L^{-2/3}\), slower than the
\(1/L\) of a free particle in a box: the zero-point mechanism of
\href{action-floor-yang-mills-gap.md}{G08} gives a softer confinement
than a hard wall.

\hypertarget{the-lattice-route-in-d3-and-the-strong-coupling-end}{%
\subsection{\texorpdfstring{3. The lattice route in \(d=3\) and the
strong-coupling
end}{3. The lattice route in d=3 and the strong-coupling end}}\label{the-lattice-route-in-d3-and-the-strong-coupling-end}}

On the periodic square lattice \((a\mathbb Z/N_sa\mathbb Z)^2\) the
Kogut--Susskind Hamiltonian is
\[H=\frac{\hbar c}{a}\Big[\frac{g_{\rm lat}^2}{2}\sum_\ell(-\Delta_\ell)
+\frac{2}{g_{\rm lat}^2}\sum_p(N-\operatorname{Re}\operatorname{tr}U_p)\Big],
\qquad g_{\rm lat}^2=g^2\hbar a=\frac a{\ell_g},\] with \(g_{\rm lat}\)
dimensionless. Theorem 1 of \href{mass-gap-obligations-lattice.md}{the
obligations map} holds verbatim in two dimensions:
\(\Delta_{a,L}=(\hbar c/a)\,\delta(g_{\rm lat};N_s,G)>0\) with a unique
physical ground state. The strong-coupling theorem holds with \(\nu=2\),
\(\Lambda_0=\{0,1\}^2\), two links and one plaquette per site:
\[\beta=\frac{4}{g_{\rm lat}^4C_2}\cdot2N=\frac{16N^2}{g_{\rm lat}^4(N^2-1)},\qquad
\Delta_{a,L}\ge\gamma_2\,\frac{g_{\rm lat}^2}{2}\,C_2\,\frac{\hbar c}{a},\]
for
\(g_{\rm lat}\ge g_0^{(2)}(N)=\big(16N^2/[(N^2-1)\beta_*(2,\{0,1\}^2)]\big)^{1/4}\),
uniformly in \(N_s\), by the same hypothesis check.

The continuum limit is \(a\to0\) at fixed \(g^2\), \(L\), that is
\(g_{\rm lat}^2\to0\) with \(N_s=L/a\) and
\(x=g^2\hbar L=g_{\rm lat}^2N_s\) fixed. The lattice gap in units of
\(\varepsilon_3\) is
\[\frac{\Delta_{a,L}}{\varepsilon_3}=\frac{\hbar c}{a\,g^2\hbar^2c}\,\delta
=\frac{\delta(g_{\rm lat};N_s,G)}{g_{\rm lat}^2},\] so the continuum
finite-volume gap exists iff
\(\delta(g_{\rm lat};x/g_{\rm lat}^2,G)/g_{\rm lat}^2\to f(x)\) as
\(g_{\rm lat}\to0\), and the infinite-volume mass is
\[\frac{mc^2}{\varepsilon_3}=\lim_{g_{\rm lat}\to0}\ \frac{\delta_\infty(g_{\rm lat})}{g_{\rm lat}^2},
\qquad\delta_\infty(g_{\rm lat})=\liminf_{N_s\to\infty}\delta(g_{\rm lat};N_s,G).\]
At strong coupling the same ratio is bounded below by \(\gamma_2C_2/2\).
The \(d=3\) conjecture on the lattice route is therefore:

\begin{quote}
\textbf{T3\(_{(3)}\).} The function
\(c(g_{\rm lat})=\delta_\infty(g_{\rm lat})/g_{\rm lat}^2\), which
satisfies \(c\ge\gamma_2C_2/2\) for \(g_{\rm lat}\ge g_0^{(2)}\), stays
bounded below by a positive constant for all \(g_{\rm lat}>0\) and has a
finite limit \(C\) as \(g_{\rm lat}\to0\).
\end{quote}

\textbf{The abelian contrast in \(d=3\).} For compact \(U(1)\) in three
dimensions Göpfert and Mack (abstract, B78) prove a gap at every
coupling, so T2\('\) holds for the abelian group here, and their Debye
mass \(m_D^2=(2\beta/a^3)e^{-\beta v(0)/2}\),
\(\beta=4\pi^2/e^2_{\rm lat}\), makes
\(c(g_{\rm lat})=\delta_\infty/g_{\rm lat}^2\) vanish with an essential
singularity as \(g_{\rm lat}\to0\): the continuum photon is massless. In
\(d=3\) the abelian/non-abelian distinction is therefore the limit of
the ratio, zero against a positive \(C\), with positivity at each
coupling shared by both. Everything in this note up to here is a
collection of known facts in one normalization; the only statement of
this programme's own is the reduction to \(f\) and its two ends.

No exponential enters. In \(d=4\) the corresponding ratio
\(\delta_\infty(g)/g^2\) is bounded below at strong coupling by the same
theorem and must vanish like \(g^{-2}e^{-1/(2b_0g^2)}\) at weak
coupling; the two ends of the \(d=4\) problem have different functional
forms and the lattice gap must interpolate between them, which is what
makes T3 in \(d=4\) a transmutation statement. In \(d=3\) both ends are
linear in \(g_{\rm lat}^2\).

\hypertarget{what-is-in-hand-and-what-is-not}{%
\subsection{4. What is in hand and what is
not}\label{what-is-in-hand-and-what-is-not}}

\begin{longtable}[]{@{}
  >{\raggedright\arraybackslash}p{(\columnwidth - 4\tabcolsep) * \real{0.3333}}
  >{\raggedright\arraybackslash}p{(\columnwidth - 4\tabcolsep) * \real{0.3333}}
  >{\raggedright\arraybackslash}p{(\columnwidth - 4\tabcolsep) * \real{0.3333}}@{}}
\toprule\noalign{}
\begin{minipage}[b]{\linewidth}\raggedright
\end{minipage} & \begin{minipage}[b]{\linewidth}\raggedright
\(d=3\)
\end{minipage} & \begin{minipage}[b]{\linewidth}\raggedright
\(d=4\)
\end{minipage} \\
\midrule\noalign{}
\endhead
\bottomrule\noalign{}
\endlastfoot
scale of the gap & \(\varepsilon_3=g^2\hbar^2c\), classical constant &
\(\hbar c\Lambda\), generated by transmutation \\
finite-volume gap & \(\varepsilon_3f(x)\), \(x=g^2\hbar L\) &
\((\hbar c/L)\,z(L)\) with \(g(L)\) running \\
small-volume end & \(f(x)=\delta_1^{(2)}x^{-2/3}[1+o(1)]\) &
\(z=\delta_1^{(3)}g(L)^{2/3}[1+O(g^{2/3})]\) \\
strong-coupling lattice end &
\(\delta_\infty/g_{\rm lat}^2\ge\gamma_2C_2/2\) &
\(\delta_\infty/g^2\ge\gamma_3C_2/2\) \\
conjecture & \(\inf_{g_{\rm lat}}c(g_{\rm lat})>0\), \(c\to C\) &
\(\delta_\infty(g)/(a\Lambda_{\rm lat}(g))\to m/(\hbar c\Lambda_{\rm lat})\) \\
expected constant & \(C=N/(2\pi)\) (Karabali--Nair, physics) & none in
closed form \\
existence & finite volume by stochastic quantisation (Chevyrev review,
passage-level companion in docs) & finite-volume UV stability only \\
\end{longtable}

Open in \(d=3\): (i) the infinite-volume limit of the finite-volume
construction; (ii) any lower bound on \(f\) away from the two ends;
(iii) the identification of \(f\)'s small-\(x\) correction with the
stochastic construction's estimates. Item (ii) is the conjecture itself,
and the one-variable form means that a single inequality
\(f(x)\ge f_->0\) for \(x\ge x_0\), proved for the continuum
finite-volume theory, together with the existence of \(\lim f\), would
settle the \(d=3\) mass gap.

\hypertarget{consequence-for-state}{%
\subsection{5. Consequence for STATE}\label{consequence-for-state}}

The \(d=3\) target is now a statement about one function of one variable
with both ends known: \(f(x)\sim\delta_1^{(2)}x^{-2/3}\) at \(x\to0\)
and \(f\to C\) conjectured at \(x\to\infty\), with the lattice ratio
\(\delta_\infty/g_{\rm lat}^2\) bounded below at strong coupling. It
replaces the \(d=3\) line of STATE. The next theorem-sized step in
\(d=3\) is a lower bound \(f(x)\ge f_->0\) on an interval
\(x\in[x_0,x_1]\) beyond the small-volume expansion, for which the
candidate tool is a variational or operator-inequality argument on the
torus Hamiltonian that keeps the zero-point mechanism of the constant
modes while controlling the nonzero modes by their Gaussian part, which
is what the weak-coupling lattice bound of STATE step 3 is meant to
prepare.
\end{document}
